% Generated draft: requires mathematical and source review.
\documentclass[11pt]{article}
\usepackage[margin=25mm]{geometry}
\usepackage{fontspec}
\setmainfont{Latin Modern Roman}
\usepackage{amsmath,amssymb,mathtools,unicode-math}
\setmathfont{Latin Modern Math}
\usepackage{enumitem,needspace,xurl,hyperref}
\setlength{\parindent}{0pt}
\setlength{\parskip}{5pt}
\setlength{\emergencystretch}{3em}
\newcommand{\N}{\mathbb{N}}
\newcommand{\Z}{\mathbb{Z}}
\newcommand{\Q}{\mathbb{Q}}
\newcommand{\R}{\mathbb{R}}
\newcommand{\C}{\mathbb{C}}
\newcommand{\F}{\mathbb{F}}
\newcommand{\A}{\mathbb{A}}
\newcommand{\PP}{\mathbb{P}}
\newcommand{\E}{\mathbb{E}}
\newcommand{\eps}{\varepsilon}
\newcommand{\dd}{\,\mathrm{d}}
\begin{document}
\section*{UnsolvedMath ID: 30004361}
\textbf{Source problem code:} catalogue entry
\subsection{Definitions and mathematical statement}
A hypergraph \(H=(V,E)\) consists of a vertex set \(V\) and a family \(E\) of subsets of \(V\), called hyperedges. It is \emph{6-uniform} if \(|e|=6\) for every \(e\in E\), and \emph{6-regular} if every vertex belongs to exactly six hyperedges. A \emph{3-coloring} is a map \(c:V\to\{1,2,3\}\); it need not use all three colors.

The supplied problem asks: does every 6-regular, 6-uniform hypergraph \(H=(V,E)\) admit a 3-coloring \(c\) such that
\[
\forall e\in E\;\forall i\in\{1,2,3\},\qquad
\bigl|\{v\in e:c(v)=i\}\bigr|\leq 3?
\]
Thus, no hyperedge may contain four or more vertices of any one color. The source does not explicitly specify whether only finite hypergraphs are intended or whether repeated hyperedges are allowed. These conventions remain unresolved; if repeated hyperedges are admitted, the degree convention must specify whether their occurrences are counted separately.
\subsection{Short English statement}
Can the vertices of every hypergraph with six vertices in each edge and six edges through each vertex be colored with three colors so that each edge has at most three vertices of any one color? The supplied record labels this problem open; its finiteness and repeated-edge conventions require clarification.
\subsection{Sources}
\begin{itemize}
\item Local UnsolvedMath registry: supplied complete record for permanent ID \(30004361\), titled \emph{Balanced Three-Colorings of Six-Uniform Hypergraphs}, in the Oberwolfach Reports Open Problems collection; exact source problem code \texttt{catalogue entry}. The record labels the statement exact and the problem open, and lists proposed year 2020. No individual proposer is supplied.
\item Supplied report citation: \emph{Combinatorics (2021)}, Oberwolfach Reports, DOI 10.4171/owr/2020/1, pp.\ 6--89. \url{https://doi.org/10.4171/owr/2020/1}. The supplied background describes these reports as documenting mathematical workshops and their open-problem discussions.
\item Supplied literature reference: Anastos--Lamaison--Steiner--Szab\'o, \url{https://www.combinatorics.org/ojs/index.php/eljc/article/view/v28i2p31}.
\item Literature assessment checked at \texttt{2026-08-22}: the supplied curation states that the cited sparse-digraph majority-coloring paper explicitly leaves the six-regular, six-uniform hypergraph coloring problem open and that no later proof was located. This is supplied dated curation derived from the Oberwolfach Reports record, not an independently verified assessment. The record requests independent verification of time-sensitive status and cited claims.
\end{itemize}
\end{document}
